\section{Problem Formulation}
\label{sec:formulation}

% New file 2026-09-15, hosting docs/PAPER.md §3 (four steps) from
% docs/problem.md workload model and resource frontier.

This section fixes the objects and the objective that design and evaluation
share, so that both use the same verdict.

\paragraph{Sessions and events.}
We model the workload as \emph{periodic interaction sessions}: updates become
eligible on a describable period, and adjacent updates depend on the same
growing history. With a common period $T$ and reference time $t_0$, the $k$-th
update of session $i$ has release time $r(i,k) = t_0 + \phi_i + kT$, where
$\phi_i \in [0, T)$ is the session's release offset. Release means eligibility:
the input is available, not that the request has entered the engine. We
distinguish four events per update: release $r(i,k)$, actual submission
$a(i,k)$, the start $s(i,k)$ of computation that depends on the target KV, and
the completion event $c(i,k)$. The rhythm must come from the workload itself;
a service-stack entrance may choose phase alignment within the application's
contract, in which case $a$ coincides with $r$ in steady state and each
session pays at most one period of alignment wait at its first submission,
reported separately. Waits introduced by deferring ready input count toward
latency, and changes to input chunking must be measured against original input
time.

\paragraph{Service objective.}
The soft real-time objective is $c(i,k) - r(i,k) \le D_i$. The completion
event is the \emph{generation completion} of the update; user-visible delivery
is modeled as a constant offset thereafter, an assumption we state explicitly.
The latency baseline is the actual submission $a(i,k)$ and must be read
together with an update-backlog metric, so that deferred submissions cannot
hide overload. Per-update generation is parameterized as $m(i,k) \le M$. For
sessions with real-time output consumption at rate $R_{\mathrm{play}}$,
bounded buffering and freshness limit how far cumulative generation may lead
consumption: over any $n$ consecutive periods, retained-and-delivered output
obeys $\sum_{j=k}^{k+n-1} N_{\mathrm{out}}(j) \le R_{\mathrm{play}} \cdot nT +
B_{\mathrm{out}}$ for a bounded lead allowance $B_{\mathrm{out}}$. A single
update may emit more than one period of content when prebuffering is allowed;
the contract bounds sustained lead, not each period in isolation.
\placeholder{Concrete target value $D$, permitted violation ratio
(parameterized sweep), observation horizon, and session-failure rule; pending
protocol freeze (measurement semantics owner).}

\paragraph{KV idle intervals.}
Let $q(i,k)$ be the last use of the target KV by update $k$. If the next use
$s(i,k+1)$ starts later and no other execution touches the state, the span
between them is a KV idle interval---the time window in which every
residency-reducing action must fit. For sessions whose bounded update
completes early and whose next update is gated by input and clock, the
interval recurs every period. Periodicity supplies the prediction; queueing,
long computations, backlogged execution, or overlapping access can shorten or
eliminate the interval. Release offsets change how sessions overlap; they
create neither idle time nor bandwidth for any single session.

\paragraph{Resource constraints.}
Physical capacity bounds allocation at every instant: with $G_{KV}$ the usable
KV pool and $b_p$ the size of allocated block $p$, $\sum_{p} b_p \le G_{KV}$,
where prefetch targets count from allocation, before their content arrives.
Peak allocation therefore depends on when eviction and restoration are
scheduled, and allocatable space does not imply content readiness. Treating
the transfer link as a renewable time resource yields a stronger candidate
feasibility condition than aggregate bandwidth checks: each session's periodic
restoration demand should receive enough link time before its deadline. This
condition is a modeling assumption to be calibrated, not a established law:
link time is not exclusively assigned---backup, gap recomputation, and
competing sessions share it---and the window from initiation to completion
includes scheduling waits. One offset assignment thus simultaneously shapes
three coupled resources: link time, the duration of GPU space occupancy, and
batch composition. The same link view bounds the total reclaimable residency
per period by the smaller of the link's per-period byte budget and the total
evictable state, which suggests---again as an assumption to verify---that the
achievable capacity extension is a property of link capability and timing
structure rather than of the per-session retention policy.
\placeholder{Batch-dependent cost function, calibrated transfer service curve,
and the feasibility condition itself; validated by independent calibration and
held-out prediction (Q5).}
